\documentclass[10pt,twocolumn,a4paper]{IEEEtran}
\usepackage[utf8]{inputenc}
\usepackage{graphicx}
\usepackage{booktabs}
\usepackage{amsmath,amssymb}
\usepackage{url}
\usepackage{hyperref}

\addtolength{\topmargin}{-0.35in}
\addtolength{\textheight}{0.7in}
\addtolength{\oddsidemargin}{-0.15in}
\addtolength{\evensidemargin}{-0.15in}
\addtolength{\textwidth}{0.3in}

\setlength{\textfloatsep}{4pt plus 1pt minus 1pt}
\setlength{\floatsep}{4pt plus 1pt minus 1pt}
\setlength{\intextsep}{4pt plus 1pt minus 1pt}
\setlength{\abovecaptionskip}{2pt}
\setlength{\belowcaptionskip}{1pt}

\title{\vspace{-0.3cm}Task C2: Abstractive Summarisation of BBC News Articles\vspace{-0.2cm}}
\author{
    \IEEEauthorblockN{KAINAT \quad Roll No. 70176544 \quad | \quad Section C}
    \IEEEauthorblockA{
        Department of Intelligent Systems, The University of Lahore \quad | \quad AI-101 Intro to AI (Instructor: Dr. Muhammad Awais) \\
        Repository: \url{https://github.com/70176544-sketch/AI101-ML-Assignment-70176544}
    }
}

\begin{document}
\maketitle

\section{Introduction}
Abstractive news summarisation condenses lengthy articles into concise, fluent single-sentence headlines that capture core events without verbatim copying. In this task, we benchmarked an extractive Lead-1 baseline, a custom GRU encoder-decoder with Luong attention trained from scratch, and a fine-tuned pretrained Transformer (\texttt{t5-small}) on the BBC XSum corpus[cite: 14]. Fine-tuned \texttt{t5-small} achieved the strongest test performance with $25.79$ ROUGE-1, $6.51$ ROUGE-2, and $20.13$ ROUGE-L, outperforming Lead-1 by $+10.47$ ROUGE-1 points and eliminating the unknown-token collapse observed in the from-scratch GRU model.

\section{Problem Statement}
Given an unstructured news article token sequence $x = (x_1, \dots, x_m)$, the task is supervised sequence-to-sequence learning $f: x \mapsto y = (y_1, \dots, y_n)$ that generates a single-sentence abstractive summary ($n \ll m$)[cite: 14]. Model quality is evaluated using ROUGE-1, ROUGE-2, and ROUGE-L F1-scores alongside factual faithfulness[cite: 14]. Key constraints include extreme compression ($18.19\times$), high lexical novelty ($81.96\%$ novel summary bigrams), and limited CPU compute on Google Colab requiring input truncation ($128$--$200$ tokens).

\section{Dataset Used}
We used the XSum (Extreme Summarisation) dataset [1] from Hugging Face (\texttt{EdinburghNLP/xsum}), comprising $226,711$ BBC news articles paired with journalist-written one-sentence summaries[cite: 14]. Following task specifications, we sampled $20,000$ training, $1,000$ validation, and $1,000$ test articles from the official splits using a fixed seed (\texttt{SEED=42})[cite: 14]. No rows were dropped during sampling; text was lowercased and tokenized via regex word/punctuation splitting for Model 1 ($10,000$-word vocabulary) and SentencePiece subwords for Model 2.

\section{Dataset Statistics}
Table~\ref{tab:stats} details token length distributions (mean, median, 95th percentile), compression ratio, lexical novelty, and OOV rate[cite: 14]. Figure~\ref{fig:eda} illustrates the right-skewed article lengths versus tightly bounded single-sentence summaries.

\begin{table}[h]
\centering
\caption{Dataset statistics for Task C2 (XSum subset).}
\label{tab:stats}
\footnotesize
\begin{tabular}{@{}ll@{}}
\toprule
\textbf{Property} & \textbf{Value} \\ \midrule
Samples (train / val / test) & 20,000 / 1,000 / 1,000 \\
Article length (mean / med / 95th) & 442.5 / 353.0 / 1094.0 tokens \\
Summary length (mean / med / 95th) & 24.3 / 24.0 / 34.0 tokens \\
Compression ratio / Novel bigrams & $18.19\times$ / 81.96\% \\
Model 1 vocab / Test OOV rate & 10,000 words / 5.65\% \\
Missing values & 0.0\% \\ \bottomrule
\end{tabular}
\end{table}

\begin{figure}[h]
\centering
\includegraphics[width=\linewidth,height=3.2cm,keepaspectratio]{figures/c2_eda_lengths.png}
\caption{Token length distributions for XSum articles (left, clipped at 1,200) and one-sentence reference summaries (right).}
\label{fig:eda}
\end{figure}

\section{Methodology}
\subsection{Learning Setting}
This is a supervised text-to-text sequence generation task[cite: 14]. We compared three approaches: (1) \textbf{Lead-1}, an extractive baseline selecting the first article sentence[cite: 14]; (2) \textbf{Model 1}, a single-layer GRU encoder-decoder with dot-product attention trained from scratch with teacher forcing; and (3) \textbf{Model 2}, a pretrained \texttt{t5-small} Transformer fine-tuned with beam search ($k=2$) and \texttt{no\_repeat\_ngram\_size=3} to prevent phrase loops[cite: 14].

\subsection{Data Split Approach}
Official train/validation/test splits were shuffled with \texttt{seed=42} before subsampling $20,000 / 1,000 / 1,000$ records[cite: 14]. Model 1's vocabulary was built strictly on the training split to avoid leakage, and all models were evaluated once on the same $1,000$ test articles[cite: 15].

\subsection{Parameter Settings}
Experiments ran on Google Colab CPU. Table~\ref{tab:params} summarizes architectures, hyperparameters, and training times[cite: 1, 16].

\begin{table}[h]
\centering
\caption{Model architectures, hyperparameters, and runtime.}
\label{tab:params}
\footnotesize
\resizebox{\columnwidth}{!}{%
\begin{tabular}{@{}lll@{}}
\toprule
\textbf{Model} & \textbf{Configuration \& Hyperparameters} & \textbf{Train Time} \\ \midrule
Lead-1 Baseline & Extractive first sentence (no training) & $<1$ s \\
Model 1 (GRU+Attn) & Emb $64$, Hid $96$, Adam ($\text{lr}=3\text{e-}3$), batch $256$, $3$ epochs & 1779.7 s \\
Model 2 (\texttt{t5-small}) & Pretrained $60$M, AdamW ($\text{lr}=3\text{e-}4$), $2$ epochs, beam $2$, $g=3$ & 6468.2 s \\ \bottomrule
\end{tabular}%
}
\end{table}

\section{Results and Analysis}
\subsection{Results}
Table~\ref{tab:results} reports test-set ROUGE F1-scores and average generated lengths[cite: 14]. Figure~\ref{fig:loss} plots training and validation cross-entropy loss curves, while Table~\ref{tab:examples} displays qualitative outputs[cite: 14].

\begin{table}[h]
\centering
\caption{Test-set ROUGE F1 (\%) and average length ($n=1,000$).}
\label{tab:results}
\footnotesize
\begin{tabular}{@{}lcccc@{}}
\toprule
\textbf{Model} & \textbf{ROUGE-1} & \textbf{ROUGE-2} & \textbf{ROUGE-L} & \textbf{Avg. Len} \\ \midrule
Lead-1 Baseline & 15.32 & 1.44 & 11.46 & 25.0 \\
Model 1 (GRU+Attn) & 14.77 & 2.13 & 12.63 & 30.3 \\
Model 2 (\texttt{t5-small}) & \textbf{25.79} & \textbf{6.51} & \textbf{20.13} & 19.8 \\ \bottomrule
\end{tabular}
\end{table}

\begin{figure}[h]
\centering
\includegraphics[width=\linewidth,height=3.0cm,keepaspectratio]{figures/c2_loss_curves.png}
\caption{Training and validation loss curves for Model 1 (left) and fine-tuned \texttt{t5-small} (right).}
\label{fig:loss}
\end{figure}

\begin{table}[h]
\centering
\caption{Sample generated summaries on test articles.}
\label{tab:examples}
\footnotesize
\resizebox{\columnwidth}{!}{%
\begin{tabular}{@{}p{1.2cm}p{6.8cm}@{}}
\toprule
\textbf{Source} & \textbf{Text / Generated Summary} \\ \midrule
\textbf{Ex. 1 Art.} & Sarah Johnson was one of 21 women heading to Liverpool when their minibus was hit by a lorry on the M62... \\
\textbf{Ref.} & A woman who was seriously hurt in a fatal hen party motorway crash is now helping other major trauma victims... \\
\textbf{Model 1} & a \texttt{<unk>} has been charged with a \texttt{<unk>} \texttt{<unk>} in the \texttt{<unk>} \texttt{<unk>} in the \texttt{<unk>} . \\
\textbf{Model 2} & A woman who died after a bus hit by a lorry in Liverpool has been given a life sentence. \\ \midrule
\textbf{Ex. 3 Art.} & Two-year-old Lane Thomas Graves had been playing in the sand near the resort's Seven Seas Lagoon when he was dragged underwater by the creature... \\
\textbf{Ref.} & Walt Disney World has unveiled a lighthouse memorial for a young boy who was killed by an alligator... \\
\textbf{Model 1} & a \texttt{<unk>} has been charged with a \texttt{<unk>} in a \texttt{<unk>} \texttt{<unk>} in the \texttt{<unk>} . \\
\textbf{Model 2} & A young boy has been killed by a sandcastle in Florida, a Florida park has said. \\ \bottomrule
\end{tabular}%
}
\end{table}

\subsection{Analysis}
\textbf{Why Lead-1 Performs Poorly:} Lead-1 achieves only $15.32$ ROUGE-1 and $1.44$ ROUGE-2 because XSum summaries are true abstractive syntheses ($81.96\%$ novel bigrams) written to hook readers rather than repeating the opening sentence, unlike traditional news datasets (e.g., CNN/DailyMail).

\textbf{Truncation Cost, Repetition, and OOV Collapse:} Truncating articles to $128$--$200$ tokens discards key facts appearing later in $442.5$-token articles (such as the memorial unveiling in Ex.\ 3). Model 1 suffers from severe mode collapse and repetitive templates (\textit{``a \texttt{<unk>} has been charged with a \texttt{<unk>}...''}) due to its $5.65\%$ OOV rate, short CPU training ($3$ epochs), and greedy decoding without n-gram blocking. Conversely, \texttt{t5-small} uses subword tokenization (zero OOV) and \texttt{no\_repeat\_ngram\_size=3}, eliminating verbatim loops.

\textbf{Hallucination Audit (20 Summaries):} Manual inspection of the first $20$ test summaries from Model 2 revealed that \textbf{12 out of 20 ($60\%$)} contained factual hallucinations not supported by the source article—for example, stating a boy was \textit{``killed by a sandcastle''} instead of an alligator (Ex.\ 3), falsely claiming a crash victim was \textit{``given a life sentence''} (Ex.\ 1), or inverting a miscarriage-of-justice acquittal into being \textit{``jailed for a year''} (Ex.\ 7). Because XSum training targets contain external background facts absent from truncated inputs, fine-tuned models learn to fabricate plausible-sounding entities.

\section{Conclusion}
Fine-tuned \texttt{t5-small} ($25.79$ ROUGE-1) substantially outperformed both the extractive Lead-1 baseline ($15.32$) and our from-scratch GRU+Attention model ($14.77$), producing fluent summaries free of unknown tokens and n-gram repetition. However, input truncation and XSum's high novelty induced frequent entity hallucinations ($60\%$ in manual audit). Next, we would train on a GPU with full $512$-token inputs and apply factual consistency reranking (e.g., NLI entailment scoring).

\newpage
\section*{AI Use Disclosure}
I used Google Gemini (November 2026) for structuring the PyTorch GRU+Attention and Hugging Face \texttt{t5-small} training scripts and formatting the LaTeX report. The following parts were generated or substantially assisted by AI: notebook training/evaluation boilerplate and LaTeX layout[cite: 18]. I verified this content by executing all cells on Google Colab, inspecting all printed ROUGE scores and generated summaries, auditing 20 summaries manually, and checking library usage against Hugging Face documentation[cite: 18]. All experiments, results and analysis in this report were produced and checked by me, and I take full responsibility for this work[cite: 18].

\begin{thebibliography}{9}
\bibitem{xsum}
S. Narayan, S. B. Cohen, and M. Lapata, ``Don't Give Me the Details, Just the Summary! Topic-Aware Convolutional Neural Networks for Extreme Summarization,'' in \textit{Proc. EMNLP}, pp. 1797--1807, 2018. \url{https://huggingface.co/datasets/EdinburghNLP/xsum}.
\bibitem{t5}
C. Raffel et al., ``Exploring the Limits of Transfer Learning with a Unified Text-to-Text Transformer,'' \textit{JMLR}, vol. 21, no. 140, pp. 1--67, 2020.
\end{thebibliography}

\end{document}
